\documentclass[11pt]{article}
\usepackage[a4paper,margin=2.4cm]{geometry}
\usepackage{amsmath,amssymb,amsthm}
\usepackage{fontspec}
\setmainfont{Noto Serif CJK KR}[ItalicFont={Noto Serif CJK KR},ItalicFeatures={FakeSlant=0.2},BoldItalicFont={Noto Serif CJK KR Bold},BoldItalicFeatures={FakeSlant=0.2}]
\setmonofont{Noto Sans Mono CJK KR}
\newfontfamily\latinfont{Latin Modern Roman}
\newcommand{\Erdos}{{\latinfont Erd\H{o}s}}
\setsansfont{Noto Sans CJK KR}
\XeTeXlinebreaklocale "ko"
\XeTeXlinebreakskip = 0pt plus 1pt
\usepackage{booktabs}
\usepackage{longtable,array,calc}
\providecommand{\tightlist}{\setlength{\itemsep}{0pt}\setlength{\parskip}{0pt}}
\usepackage{enumitem}
\usepackage[hidelinks]{hyperref}
\linespread{1.25}
\setlength{\parindent}{0pt}
\setlength{\parskip}{0.5em}
\newtheorem{theorem}{정리}
\newtheorem{lemma}[theorem]{보조정리}
\theoremstyle{definition}
\newtheorem{definition}[theorem]{정의}
\newtheorem{example}[theorem]{예}
\DeclareMathOperator{\lcm}{lcm}
\DeclareMathOperator{\spf}{spf}
\DeclareMathOperator{\spn}{span}
\renewcommand{\contentsname}{목차}

\title{학습 노트}
\author{hoon (regx64)}
\date{2026년 10월 12일 기준}
\begin{document}
\maketitle
\tableofcontents

%======================================================
\section{미적분 → 다변수 미적분}
목표는 비압축성 나비에–스토크스 방정식을 수학적으로 읽는 것이다.
\[
\rho\Big(\frac{\partial \mathbf u}{\partial t}+(\mathbf u\cdot\nabla)\mathbf u\Big)=-\nabla p+\mu\nabla^2\mathbf u+\mathbf f,\qquad \nabla\cdot\mathbf u=0.
\]
로드맵: 1변수 미적분 → 다변수 미적분 → 벡터 미적분 → ODE/PDE → 선형대수.

\subsection{1변수}
\begin{itemize}[nosep]
\item 극한: 대입하면 $0/0$이 되는 경우 인수분해로 처리, 좌극한·우극한.
\item 연속: $\lim_{x\to a}f(x)=f(a)$. 좌극한$=$우극한은 극한이 \emph{존재}한다는 뜻일 뿐이다. 연속이려면 $f(a)$가 정의되어 있고 그 값과 같아야 한다.
\item 미분계수: $f'(a)=\lim_{h\to0}\dfrac{f(a+h)-f(a)}{h}$. 할선 기울기가 접선 기울기로 수렴한다.
\item 도함수 규칙: 거듭제곱, 곱의 법칙, 연쇄법칙 $\big(f(g(x))\big)'=f'(g(x))\,g'(x)$.
\item 활용: 증가·감소, 극값, 부호표.
\item 부정적분, 정적분, 미적분학의 기본정리 $\dfrac{d}{dx}\displaystyle\int_a^x f(t)\,dt=f(x)$.
\end{itemize}

\subsection{다변수}
\begin{itemize}[nosep]
\item 편미분, 다변수 연쇄법칙.
\item 이계편미분, Clairaut 정리 $f_{xy}=f_{yx}$ (2계 편도함수가 연속일 때).
\item 그래디언트 $\nabla f$: 등고선에 수직이고 가장 가파른 오르막 방향.
\item 발산 $\nabla\cdot\mathbf F=\partial_xF_1+\partial_yF_2+\partial_zF_3$.
\item 회전 $\nabla\times\mathbf F=(\partial_yF_3-\partial_zF_2,\ \partial_zF_1-\partial_xF_3,\ \partial_xF_2-\partial_yF_1)$.
\item 다변수 극값: 임계점, 안장점, 이차미분판정법 $D=f_{xx}f_{yy}-f_{xy}^2$.
\end{itemize}

\begin{center}
\begin{tabular}{ll}
\toprule
항 & 수학\\
\midrule
$-\nabla p$ & 압력의 그래디언트\\
$\mu\nabla^2\mathbf u$ & 라플라시안 (점성)\\
$\nabla\cdot\mathbf u=0$ & 발산 0 (비압축)\\
$(\mathbf u\cdot\nabla)\mathbf u$ & 대류항, 연쇄법칙에서 나옴\\
\bottomrule
\end{tabular}
\end{center}

자주 한 실수: 연쇄법칙에서 안쪽 함수의 미분 누락, 회전 공식의 항 순서 반대로 적용, 연립방정식 대입 계산 실수.

다음: $\nabla\times(\nabla f)=\mathbf 0$, $\nabla\cdot(\nabla\times\mathbf F)=0$, $\nabla\cdot(\nabla f)=\nabla^2 f$ 같은 관계를 묻는 종합 문제.

%======================================================
\section{푸리에 변환}
2026년 10월 10일 시작. 선수지식까지 포함해 엄밀하게 배우고 자유롭게 적용하는 수준이 목표다. 교재는 한국어 \texttt{materials/fourier/fourier-transform.pdf} (103쪽).

\subsection{교재 구성}
\begin{enumerate}[nosep,start=0]
\item 준비: 삼각함수 (라디안, 단위원, 덧셈정리, 작은 각 근사)
\item 미적분 도구와 복소지수 $e^{j\theta}=\cos\theta+j\sin\theta$
\item 선형대수: 기저와 좌표, 행렬, 내적공간, 정규직교기저, 고유벡터
\item 이산 푸리에 변환
\item 푸리에 급수
\item 푸리에 변환과 성질, 변환쌍 표
\item 델타 함수
\item LTI 시스템과 합성곱
\item 표본화 정리, FFT
\item 불확정성 원리
\item 확장: STFT, 2차원 변환, 라플라스 변환, 열방정식
\end{enumerate}
규약은 전기공학식이다.
\[
X(\omega)=\int_{-\infty}^{\infty}x(t)\,e^{-j\omega t}\,dt,\qquad
x(t)=\frac{1}{2\pi}\int_{-\infty}^{\infty}X(\omega)\,e^{j\omega t}\,d\omega.
\]
핵심 관점: 푸리에 변환은 ``시각별 값''이라는 좌표를 ``주파수별 성분''이라는 좌표로 바꾸는 기저 변환이다.

\subsection{진도}
0장 진행 중. 반지름 $c$인 원 위의 점은 $(c\cos\theta,\ c\sin\theta)$이고, 단위원은 $c=1$인 경우다.

\begin{center}
\begin{tabular}{cc}
\toprule
연습문제 & 필요한 절\\
\midrule
0.1 & 0.1\\
0.2 & 0.3, 0.6\\
0.3 & 0.2\\
0.4 & 0.4\\
0.5 & 0.5\\
0.6, 0.7, 0.9 & 0.7\\
0.8 & 0.6\\
\bottomrule
\end{tabular}
\end{center}
속도 기준은 하루 1–2절이다. 증명을 다시 쓰고 연습문제를 보지 않고 풀 수 있으면 완료로 본다.

%======================================================
\section{선형대수}
2026년 10월 9일 시작. 기한 없이 전 분야를 다룬다. 시작 시점에는 연립방정식은 풀 수 있지만 행렬로 푸는 법은 몰랐다.

\begin{itemize}[nosep]
\item 선형결합, $\spn$, 일차독립, 기저와 차원. 기저가 정해지면 좌표는 유일하다.
\item 행렬–벡터 곱은 열들의 선형결합: $A\mathbf x=x_1\mathbf a_1+\cdots+x_n\mathbf a_n$.
\item 연립방정식 $c_1+c_2=5,\ c_1-c_2=1$은 한 줄로
$\begin{pmatrix}1&1\\1&-1\end{pmatrix}\begin{pmatrix}3\\2\end{pmatrix}=\begin{pmatrix}5\\1\end{pmatrix}$.
\item 선형변환 $L(a\mathbf x+b\mathbf y)=aL(\mathbf x)+bL(\mathbf y)$. $AB$는 ``$B$ 다음 $A$''이고 일반적으로 $AB\ne BA$.
\item 역행렬이 있으면 $A\mathbf x=\mathbf b\Rightarrow\mathbf x=A^{-1}\mathbf b$.
\end{itemize}
\subsection{응용: 페이지랭크 (2026-08-31)}

페이지랭크는 행렬 계산이 정의 자체다.

\begin{itemize}
\tightlist
\item
  전이행렬 \(M\): 페이지 \(j\)가 링크 \(L_j\)개를 가지면
  \(M_{ij}=1/L_j\) (\(j\to i\) 링크가 있을 때). 열의 합이 1.
\item
  감쇠계수 \(d\) (보통 0.85)를 넣은 구글 행렬
\end{itemize}

\[
G=dM+\frac{1-d}{n}\mathbf 1\mathbf 1^{\top}
\]

\begin{itemize}
\tightlist
\item
  페이지랭크 벡터 \(\mathbf r\)은 고유값 1에 대한 고유벡터:
  \(G\mathbf r=\mathbf r\), \(\sum r_i=1\)
\item
  계산은 거듭제곱법:
  \(\mathbf r_{t+1}=dM\mathbf r_t+\frac{1-d}{n}\mathbf 1\)을 수렴할
  때까지 반복
\item
  실제 웹 규모에서는 \(M\)을 희소행렬로 저장
\end{itemize}

예: 링크 \(1\to2,\ 1\to3,\ 2\to3,\ 3\to1\), \(d=0.85\)이면
\(\mathbf r\approx(0.388,\ 0.215,\ 0.397)\). \(G\)의 나머지 고유값은
\(-0.425\pm0.425i\) (크기 약 0.60)라서 반복할 때마다 오차가 약 0.6배로
줄어든다.


다음: 가우스 소거법, 행렬식, 랭크, 고유값, 내적공간.

%======================================================
\section{P vs NP}
\begin{itemize}[nosep]
\item P: 다항 시간 $O(n^c)$에 풀 수 있는 문제. NP: 답이 주어지면 다항 시간에 검증할 수 있는 문제.
\item 부분집합 합: 부분집합 $2^n$개를 다 시도하면 지수 시간, 주어진 부분집합의 합 확인은 $O(n)$.
\item ``어려워 보인다''와 ``어렵다고 증명됐다''는 다르다. 소수 판정은 2002년 AKS로 P임이 증명됐다.
\item $n=100$이면 $n^3=10^6$, $2^{100}\approx1.27\times10^{30}$. 초당 $10^9$회 연산이면 $2^{100}$번은 약 $4\times10^{13}$년.
\end{itemize}
형식적 정의: $L\in\mathrm{NP}$ $\iff$ 다항식 $p$와 다항 시간 검증기 $V$가 있어
\[
x\in L\iff \exists\,w,\ |w|\le p(|x|):\ V(x,w)=1.
\]
다음: 다항 시간 환원, NP-완전, Cook–Levin 정리.

%======================================================
\section{도형의 방정식}
\subsection{직접 유도한 접선 공식}
원 $x^2+y^2=r^2$ 밖의 점 $P(p,q)$에서 그은 접선. $d^2=p^2+q^2$, $\ell=\sqrt{d^2-r^2}$일 때
\[
r(px+qy)\pm\ell(py-qx)=rd^2.
\]
$P$ 중심 꼴은 $r(uX+vY)\pm\ell(uY-vX)=0$, 복소수 꼴은 법선 $N=(u+vi)(r\pm\ell i)$에 대해 $N_1(x-p)+N_2(y-q)=0$이다. Lean 4로 임의의 가환환 위에서 검증했다.

\subsection{삼각형의 내심}
네 가지 동치 표현: 각의 이등분선의 교점, 세 변까지 거리가 같은 점, $\angle BIC=90^\circ+\tfrac A2$, $I=\dfrac{aA+bB+cC}{a+b+c}$.

세 직선이 주어지면 점–직선 거리를 같게 놓고, 부호는 맞은편 꼭짓점을 대입해 정한다.
\begin{itemize}[nosep]
\item $x=0,\ y=0,\ 3x+4y-12=0$ → $I(1,1)$, $r=1$.
\item $y=0,\ 5x-12y+60=0,\ 3x+4y-12=0$ → $I(-2,2)$, $r=2$.
\end{itemize}

\subsection{외심 빠르게 구하기 (2026-09-29)}

\begin{enumerate}
\def\labelenumi{\arabic{enumi}.}
\tightlist
\item
  특수한 꼴부터 확인

  \begin{itemize}
  \tightlist
  \item
    직각삼각형: 빗변의 중점
  \item
    정삼각형: 무게중심과 같음
  \item
    이등변삼각형: 꼭지각의 이등분선(밑변의 수직이등분선) 위
  \item
    직사각형·정다각형: 대각선의 교점 / 중심
  \end{itemize}
\item
  일반: 한 꼭짓점을 원점으로 평행이동하면 거리 제곱 조건
  \(|P-A|^2=|P-B|^2\)에서 이차항이 지워져 일차식만 남는다. \(A\)를
  원점으로 옮긴 뒤 \(B'=(b_1,b_2)\), \(C'=(c_1,c_2)\)라 하면
\end{enumerate}

\[
D=2(b_1c_2-b_2c_1),\qquad U_x=\frac{c_2(b_1^2+b_2^2)-b_2(c_1^2+c_2^2)}{D},\qquad U_y=\frac{b_1(c_1^2+c_2^2)-c_1(b_1^2+b_2^2)}{D}
\]

외심은 \(A+(U_x,U_y)\). 예: \(A(1,1),B(5,1),C(1,4)\) → \((3,\tfrac52)\)
(직각삼각형이라 빗변 중점과 일치). 3. 반지름만 필요하면

\[
R=\frac{a}{2\sin A}=\frac{abc}{4S}
\]

\begin{enumerate}
\def\labelenumi{\arabic{enumi}.}
\setcounter{enumi}{3}
\tightlist
\item
  사각형 이상은 원에 내접할 때(대각의 합이 \(180^\circ\))만 외심이 있고,
  그때는 아무 세 꼭짓점의 외심이 전체의 외심이다.
\end{enumerate}

\subsection{기본 공식 정리 (교과서 범위)}

\begin{itemize}
\tightlist
\item
  원과 직선: 중심--직선 거리 \(d\)와 \(r\) 비교 (\(d<r\) 두 점, \(d=r\)
  접함, \(d>r\) 만나지 않음). 원의 식이 복잡하면 판별식보다 거리 공식이
  빠르다.
\item
  현의 길이 \(2\sqrt{r^2-d^2}\). 원 위의 점과 직선 사이 거리는 최대
  \(d+r\), 최소 \(d-r\) (\(d>r\)).
\item
  접선

  \begin{itemize}
  \tightlist
  \item
    원 위의 점 \((x_1,y_1)\): \(x_1x+y_1y=r^2\), 중심 \((a,b)\)이면
    \((x_1-a)(x-a)+(y_1-b)(y-b)=r^2\)
  \item
    기울기 \(m\): \(y=mx\pm r\sqrt{m^2+1}\)
  \item
    원 밖의 점: \(y-y_1=m(x-x_1)\)로 놓고 거리 \(=r\) (수직 접선 누락
    주의)
  \end{itemize}
\item
  접선의 길이: 일반형 \(x^2+y^2+Ax+By+C=0\)에 점을 대입한 값의 제곱근
\end{itemize}

\[
\ell=\sqrt{x_1^2+y_1^2+Ax_1+By_1+C}
\]

\begin{itemize}
\tightlist
\item
  극선(두 접점을 잇는 직선): \(x_1x+y_1y=r^2\)
\item
  두 원의 위치 관계: 중심 거리 \(d\)와 \(r_1\pm r_2\) 비교
\item
  공통현: 두 원의 식을 뺀다. 두 교점을 지나는 원: \(C_1+kC_2=0\)
  (\(k\ne-1\)) 또는 \(C_1+k\cdot(\text{공통현})=0\)
\item
  공통접선 길이: 외접선 \(\sqrt{d^2-(r_1-r_2)^2}\), 내접선
  \(\sqrt{d^2-(r_1+r_2)^2}\)
\item
  아폴로니우스의 원: \(PA:PB=m:n\) (\(m\ne n\))의 자취. 내분점과
  외분점이 지름의 양 끝
\end{itemize}

\subsection{시험 문제: 원다발 + 근과 계수의 관계}

원 \(x^2+y^2=5\)와 다른 한 원의 공통현이 \(3x+2y-5=0\)일 때, 두 교점을
지나고 \(y\)축에 접하는 원들의 반지름의 곱은?

\begin{enumerate}
\def\labelenumi{\arabic{enumi}.}
\tightlist
\item
  교점을 지나는 원: \(x^2+y^2-5+k(3x+2y-5)=0\) (원 + \(k\cdot\)직선이 원
  두 개를 섞는 것보다 가볍다)
\item
  중심 \(\left(-\tfrac{3k}{2},-k\right)\)
\item
  \(y\)축에 접함 ⇔ \(x=0\)을 넣은 \(y^2+2ky-5-5k=0\)이 중근 ⇔
  \(k^2+5k+5=0\)
\item
  반지름 \(=|\text{중심의 }x\text{좌표}|=\tfrac32|k|\)
\item
  \(k\)를 직접 구하지 않고 \(k_1k_2=5\)만 쓴다:
\end{enumerate}

\[
\frac94|k_1k_2|=\frac94\cdot5=\frac{45}{4}
\]

판별식 \(25-20>0\)이라 원이 실제로 두 개다.

\subsection{원본 파일 (이 저장소에 없음)}

``Unknown purpose inquiry'' 채팅(2026-09-27)에서 만든 다음 파일은 다시
올려야 들어온다: 기본·심화 공식 정리 PDF, 공식별 예제 풀이 PDF, Lean 4
검증 파일 2개(접선 공식, 회전·복소수 해석), 접선 공식 유도 애니메이션
MP4.


%======================================================
\section{이산수학 (논리)}

시작: 2026-09-28. CS 수학 로드맵의 1단계.

\subsection{확정한 학습 순서}

\begin{enumerate}
\def\labelenumi{\arabic{enumi}.}
\tightlist
\item
  이산수학(논리, 집합, 함수, 관계, 귀납법) + 미적분 I
\item
  미적분 II\textasciitilde III (다변수, 편미분, 벡터 미적분)
\item
  선형대수
\item
  확률과 통계
\item
  이산수학 심화(점근 표기, 점화식, 그래프) + 미분방정식 + 신호와 시스템
  수학
\item
  진로별 심화 (ML: 최적화·정보이론, 이론전산: 추상대수·오토마타)
\end{enumerate}

과목마다 정의 → 정리 이해 → 손 계산 → CS/EE 응용 순서. 교재 후보는
Rosen(이산수학), Strang(선형대수).

\subsection{1강: 명제논리}

연산자 \(\neg,\ \wedge,\ \vee,\ \to,\ \leftrightarrow\). 조건문은 전제가
거짓이면 항상 참이다.

\begin{longtable}[]{@{}lll@{}}
\toprule\noalign{}
\(p\) & \(q\) & \(p\to q\) \\
\midrule\noalign{}
\endhead
\bottomrule\noalign{}
\endlastfoot
T & T & T \\
T & F & F \\
F & T & T \\
F & F & T \\
\end{longtable}

동치 법칙:

\[
\neg(p\wedge q)\equiv\neg p\vee\neg q,\qquad \neg(p\vee q)\equiv\neg p\wedge\neg q,\qquad p\to q\equiv\neg p\vee q\equiv\neg q\to\neg p
\]

역 \(q\to p\)와 이 \(\neg p\to\neg q\)는 원래 명제와 동치가 아니다.

문제와 답:

\begin{enumerate}
\def\labelenumi{\arabic{enumi}.}
\tightlist
\item
  \(\neg(p\to q)\equiv p\wedge\neg q\) --- 진리표의 두 열이 (F, T, F,
  F)로 같다.
\item
  \(\neg(p\vee(q\wedge r))\equiv\neg p\wedge(\neg q\vee\neg r)\) ---
  괄호를 지켜야 한다.
\item
  ``비가 오면 우산을 챙긴다'' → 대우: 우산을 안 챙기면 비가 오지 않는다.
  역: 우산을 챙기면 비가 온다. 이: 비가 안 오면 우산을 안 챙긴다.
\end{enumerate}

응용: AND/OR/NOT 게이트, NAND만으로 모든 회로 구성,
\texttt{!(a\ \&\&\ b)} ↔ \texttt{!a\ \textbar{}\textbar{}\ !b}.

\subsection{2강: 술어논리}

술어 \(P(x)\), 정의역, 한정자 \(\forall,\exists\). 부정 규칙:

\[
\neg\forall x\,P(x)\equiv\exists x\,\neg P(x),\qquad \neg\exists x\,P(x)\equiv\forall x\,\neg P(x)
\]

한정자 순서는 바꿀 수 없다. 실수에서 \(\forall x\,\exists y\,(x+y=0)\)은
참, \(\exists y\,\forall x\,(x+y=0)\)은 거짓.

문제:

\begin{enumerate}
\def\labelenumi{\arabic{enumi}.}
\tightlist
\item
  정의역 \(\{1,2,3,4\}\), \(P(x)\) = ``\(x>3\)'': (a)
  \(\forall x\,P(x)\) (b) \(\exists x\,P(x)\) (c)
  \(\forall x\,\neg P(x)\) (d) \(\exists x\,\neg P(x)\)
\item
  부정을 술어 바로 앞까지: (a) \(\forall x\,(x>0\to x^2>0)\) (b)
  \(\exists x\,(x\text{는 소수}\wedge x\text{는 짝수})\)
\item
  정수에서 \(Q(x,y)\) = ``\(x<y\)'': (a) \(\forall x\,\exists y\,Q\) (b)
  \(\exists y\,\forall x\,Q\)
\end{enumerate}

답: 1. F, T, F, T. 2. (a) \(\exists x\,(x>0\wedge x^2\le 0)\) (b)
\(\forall x\,(x\text{는 소수가 아님}\vee x\text{는 홀수})\). 3. (a) 참
(\(y=x+1\)) (b) 거짓 (가장 큰 정수가 없음).

응용: SQL \texttt{NOT\ EXISTS}, 프로그램 검증에서 버그 찾기 = 반례 하나
찾기.

\subsection{자주 한 실수}

\begin{itemize}
\tightlist
\item
  \(\wedge\)는 둘 다 참일 때만 참인데 하나만 참이어도 T로 적음
\item
  \(\wedge,\vee\)가 섞인 식에 드모르간을 쓸 때 괄호 누락
\end{itemize}

\subsection{도구}

문제마다 답 입력기(HTML)를 만들어 쓴다. 진리표 칸은 클릭으로 T/F 토글,
기호는 버튼 입력, 결과는 복사해서 채팅에 붙인다.


\section{집합과 경우의 수 (고1)}

2026-09-29. 오래 붙잡고 씨름하는 용도로 만든 집합 문제 두 개.

\subsection{문제 1: 합이 11이 되는 쌍 금지}

\(U=\{1,\dots,10\}\), 조건 (가): \(A\)의 서로 다른 두 원소 \(a,b\)에
대해 \(a+b\ne11\).

\begin{enumerate}
\def\labelenumi{\arabic{enumi}.}
\tightlist
\item
  조건을 만족하는 \(A\)의 개수 (공집합 포함)
\item
  조건을 만족하고 \(n(A)=4\)인 \(A\)의 개수
\item
  \(A,B\)가 모두 조건을 만족하고 \(A\cap B=\varnothing\)인 순서쌍
  \((A,B)\)의 개수
\end{enumerate}

\textbf{핵심.} \(U\)는 쌍
\(\{1,10\},\{2,9\},\{3,8\},\{4,7\},\{5,6\}\)으로 나뉘고 쌍끼리는
독립이다. 짝이 이미 고정돼 있으므로 \(10\cdot8\cdot6\cdot4\cdot2\) 같은
짝짓기 수가 아니라 \textbf{쌍마다 상태를 세서 곱한다}.

\begin{enumerate}
\def\labelenumi{\arabic{enumi}.}
\tightlist
\item
  한 쌍에서 가능한 상태는 \(\varnothing,\{a\},\{11-a\}\)의 3가지 →
  \(3^5=243\)
\item
  정확히 4개 쌍에서 하나씩, 한 쌍은 비움 → \(\binom51\cdot2^4=80\)
\item
  한 쌍의 두 원소를 각각 \(A\), \(B\), 없음에 배정(\(3^2=9\))하고 ``둘
  다 \(A\)'', ``둘 다 \(B\)''를 뺌 → 쌍당 7가지, \(7^5=16807\)
\end{enumerate}

\subsection{문제 2: 연속한 두 수 금지}

\(U=\{1,\dots,12\}\), 조건 (가): \(A\)는 연속한 두 자연수를 동시에 갖지
않는다.

\begin{enumerate}
\def\labelenumi{\arabic{enumi}.}
\tightlist
\item
  조건을 만족하는 \(A\)의 개수
\item
  조건을 만족하고 \(n(A)=5\)인 \(A\)의 개수
\item
  조건을 만족하고 3의 배수를 하나도 갖지 않는 \(A\)의 개수
\end{enumerate}

\textbf{1. 점화식.} \(\{1,\dots,n\}\)에서의 개수를 \(a_n\)이라 하자.
\(n\notin A\)이면 \(a_{n-1}\)가지, \(n\in A\)이면 \(n-1\notin A\)이므로
\(a_{n-2}\)가지다.

\[
a_n=a_{n-1}+a_{n-2},\qquad a_1=2,\ a_2=3
\]

\(2,3,5,8,13,21,34,55,89,144,233,377\) → \(a_{12}=377\) (피보나치 수).

\textbf{2. 빈칸 확보.} 원소 \(x_1<\dots<x_5\)에서 \(y_i=x_i-(i-1)\)로
바꾸면 \(1\le y_1<\dots<y_5\le 8\)인 아무 5개와 일대일 대응한다.

\[
\binom{12-5+1}{5}=\binom85=56
\]

\textbf{3. 덩어리로 쪼개기.} \(3,6,9,12\)를 빼면
\(\{1,2\},\{4,5\},\{7,8\},\{10,11\}\)의 네 덩어리가 남고, 덩어리 사이는
3의 배수가 끼어 있어 서로 영향이 없다. 두 원소 덩어리에서 가능한 선택은
\(\varnothing,\{a\},\{a+1\}\)의 3가지 → \(3^4=81\)

\subsection{검증}

모든 답을 전수 탐색으로 확인했다: 243, 80, 16807, 377, 56, 81.


\section{지수의 확장}

2026-08-23.

\subsection{유리수 지수}

\(n^{1/k}\)는 \(k\)번 곱해서 \(n\)이 되는 양수, 즉 \(\sqrt[k]{n}\)이다.

\[
n^{1/k}=\sqrt[k]{n},\qquad n^{m/k}=\big(\sqrt[k]{n}\big)^m=\sqrt[k]{n^m}\qquad(n>0)
\]

예: \(8^{1/3}=2\), \(8^{2/3}=4\).

\subsection{무리수 지수}

\(n^{\sqrt k}\)에서 \(\sqrt k\)는 (완전제곱이 아니면) 무리수라서 ``몇 번
곱한다''로는 정의할 수 없다. 두 가지 방법으로 정의한다.

\begin{enumerate}
\def\labelenumi{\arabic{enumi}.}
\tightlist
\item
  유리수 근삿값의 극한: \(\sqrt2=1.4,\ 1.41,\ 1.414,\dots\)이면
  \(2^{1.4},2^{1.41},2^{1.414},\dots\)가 한 값으로 수렴한다.
\item
  지수·로그 함수로:
\end{enumerate}

\[
n^{x}=e^{x\ln n}\qquad\Rightarrow\qquad n^{\sqrt k}=e^{\sqrt k\,\ln n}
\]

예: \(2^{\sqrt2}=e^{\sqrt2\ln2}\approx2.6651\).

참고: \(2^{\sqrt2}\)는 겔폰트--슈나이더 정리에 의해 초월수다.


\section{Shapley value (SHAP의 수학)}

2026-10-01. 설명 가능한 AI(XAI)의 SHAP 기법 밑바탕인 협력게임 이론.

\subsection{정의}

특성 집합 \(N\) (\(|N|=n\)), 부분집합 \(S\)만 알 때의 예측값을 가치
함수로 둔다. 빠진 특성은 배경 데이터 \(x'\)로 채워 평균낸다.

\[
v(S)=\mathbb E_{x'}\big[f(x_S,\ x'_{\bar S})\big],\qquad v(\varnothing)=\mathbb E_{x'}[f(x')]
\]

특성 \(i\)의 기여도:

\[
\phi_i=\sum_{S\subseteq N\setminus\{i\}}\frac{|S|!\,(n-|S|-1)!}{n!}\Big[v(S\cup\{i\})-v(S)\Big]
\]

가중치는 ``특성을 무작위 순서로 하나씩 넣을 때 \(i\) 바로 앞에 정확히
\(S\)가 들어와 있을 확률''이다.

효율성 (기여도의 합이 보존됨):

\[
\sum_{i=1}^{n}\phi_i=v(N)-v(\varnothing)
\]

\subsection{2개 특성 예 (집값: 면적 A, 역세권 B)}

\(v(\varnothing)=100,\ v(\{A\})=150,\ v(\{B\})=120,\ v(\{A,B\})=200\)

\[
\phi_A=\tfrac12(150-100)+\tfrac12(200-120)=65,\qquad \phi_B=\tfrac12(120-100)+\tfrac12(200-150)=35
\]

합 \(100=v(\{A,B\})-v(\varnothing)\).

\subsection{\texorpdfstring{연습문제 (\(n=3\))}{연습문제 (n=3)}}

\(v(\varnothing)=10,\ v(A)=20,\ v(B)=10,\ v(C)=10,\ v(AB)=30,\ v(AC)=20,\ v(BC)=10,\ v(ABC)=40\)

\(n=3\)의 가중치: \(|S|=0\)이면 \(\tfrac13\), \(|S|=1\)이면
\(\tfrac16\), \(|S|=2\)이면 \(\tfrac13\).

\[
\phi_A=\tfrac13(10)+\tfrac16(20)+\tfrac16(10)+\tfrac13(30)=\tfrac{55}{3},\qquad \phi_B=\tfrac{25}{3},\qquad \phi_C=\tfrac{10}{3}
\]

합 \(30=v(ABC)-v(\varnothing)\).

\textbf{이상한 점.} \(C\)는 혼자서도, \(A\)나 \(B\) 하나와 함께여도 값을
전혀 바꾸지 않는다. 그런데 \(A,B\)가 이미 있을 때만 \(+10\)을 만든다. 이
3자 상호작용이 셋에게 나눠지면서 \(\phi_C=\tfrac{10}{3}>0\)이 된다.
Shapley value는 상호작용을 쪼개서 개별 기여처럼 보이게 만든다. 이게
SHAP을 ``특성 중요도''로 읽을 때 생기는 비판의 출발점이다.

\subsection{실제 계산}

\(2^n\)개 부분집합을 다 볼 수 없어서 근사한다: KernelSHAP(가중
선형회귀), TreeSHAP(트리 구조로 다항 시간), DeepSHAP(DeepLIFT 기반).


\section{힐베르트 11번·24번 문제}

2026-10-03.

\subsection{11번: 대수적 수 계수 이차형식의 해}

이차형식 \(Q(x)=\sum a_{ij}x_ix_j\) (\(a_{ij}\in K\), \(K\)는 수체)의
비자명해. 답은 해를 어디서 찾느냐에 따라 갈린다.

\begin{enumerate}
\def\labelenumi{\arabic{enumi}.}
\tightlist
\item
  \textbf{대수적 폐포 \(\overline{\mathbb Q}\):} 변수가 2개 이상이면
  항상 있다. 대각화 후 \(ax^2+by^2=0\)에서 \(x=\sqrt{-b/a}\,y\).
\item
  \textbf{수체 \(K\):} 항상은 아니지만 완전히 판정된다
  (하세--민코프스키).
\end{enumerate}

\[
Q\text{가 }K\text{에서 비자명해를 가짐}\iff\text{모든 완비화 }K_v\text{에서 비자명해를 가짐}
\]

\begin{itemize}
\tightlist
\item
  반례: \(x^2+y^2+z^2=0\) (\(\mathbb R\)에서 막힘), \(x^2+y^2-3z^2=0\)
  (\(\mathbb Q_3\)에서 막힘)
\item
  변수 5개 이상이면 비아르키메데스 자리에서는 항상 해가 있어서 실수
  자리의 부호만 보면 된다
\end{itemize}

\begin{enumerate}
\def\labelenumi{\arabic{enumi}.}
\setcounter{enumi}{2}
\tightlist
\item
  \textbf{정수해 \(\mathcal O_K\):} 국소--대역 원리가 깨진다.
\end{enumerate}

\[
x^2+11y^2=3
\]

은 \(\mathbb R\)과 모든 \(\mathbb Z_p\)에서 해가 있지만 정수해가 없다.
같은 종(genus)의 \(3x^2+2xy+4y^2\)는 \((1,0)\)에서 3을 표현한다.
판별식은 둘 다 \(-44\).

그래서 11번은 보통 ``부분적으로 해결''로 분류된다. 유리해는
하세(1923--24)가 해결, 정수해는 지겔 등이 상당 부분 해결.

\subsubsection{직접 세워 본 국소--대역 원리 (종 버전)}

조건 없는 정수해 버전은 \(x^2+11y^2=3\) 때문에 불가능하다. 여기서 ``조건
없다''는 모든 \(Q\), 모든 \(c\)에 예외 없이 적용된다는 뜻이다. 그래서
가정을 강화하거나 결론을 약화해야 한다. 결론을 약화한 버전 (새 정리는
아니고 종 이론의 알려진 결과를 재구성한 것):

\textbf{설정.} \((V,Q)\)는 \(\mathbb Q\) 위의 비퇴화 이차공간,
\(L\subset V\)는 \(\mathbb Z\)-격자, \(L_p=L\otimes\mathbb Z_p\),
\(c\in\mathbb Z\setminus\{0\}\). \(\operatorname{gen}(L)\)은 모든
\(p\)에서 \(M_p\cong L_p\)인 격자 \(M\)의 집합.

\textbf{정리.}

\[
c\in Q(V\otimes\mathbb R)\ \text{이고 모든 }p\text{에서 }c\in Q(L_p)\iff\text{어떤 }M\in\operatorname{gen}(L)\text{에 대해 }c\in Q(M)
\]

\textbf{증명} (\(\Rightarrow\)).

\begin{enumerate}
\def\labelenumi{\arabic{enumi}.}
\tightlist
\item
  \(Q\perp\langle-c\rangle\)가 모든 자리에서 등방이므로
  하세--민코프스키로 \(\mathbb Q\)에서도 등방. 따라서 \(Q(v)=c\)인
  \(v\in V\)가 있다.
\item
  \(v\)의 분모에 나오는 소수는 유한 개이므로 유한집합 \(S\) 밖의
  \(p\)에서 \(v\in L_p\).
\item
  \(p\in S\)에서 가정으로 \(Q(w_p)=c\)인 \(w_p\in L_p\)가 있고, 비트
  정리로 \(\sigma_p(w_p)=v\)인 \(\sigma_p\in O(V_p)\)가 있다.
\item
  \(M_p=\sigma_p(L_p)\) (\(p\in S\)), \(M_p=L_p\) (\(p\notin S\))로 두면
  \(M=V\cap\bigcap_pM_p\)는 격자다.
\item
  모든 \(p\)에서 \(v\in M_p\)이므로 \(v\in M\)이고
  \(M\in\operatorname{gen}(L)\). \(\blacksquare\)
\end{enumerate}

\textbf{따름정리.} 종에 류(class)가 하나뿐이면 정수해 국소--대역 원리가
그대로 성립한다. \(x^2+y^2\), \(x^2+y^2+z^2\), \(x^2+y^2+z^2+w^2\)가
그렇다 (세 제곱수 정리: \(n\ne4^a(8b+7)\)).

\textbf{더 정밀한 결과.} 부정부호 3변수 이상은 스피너 종 하나에 류
하나라서 장애물이 스피너 예외로 계산된다. 정수점 장애물은 브라우어--마닌
장애물뿐(Colliot-Thélène--Xu). 정부호 5변수 이상은 충분히 큰 \(c\)에서
성립(Tartakowsky).

\subsection{24번: 증명의 단순성}

1900년 발표 목록에서 빠졌다가 2000년 뤼디거 틸레가 힐베르트의 노트에서
찾아냈다. 내용: 가장 단순한 증명을 판정하는 기준을 세우고, 증명 방법
자체의 이론을 만들 것. 참·거짓 문제가 아니라 연구 방향이라 ``해결''
판정이 없다.

진척 (방향별):

\begin{itemize}
\tightlist
\item
  \textbf{증명 길이:} 쿡--렉하우(1979)의 증명 체계. 비둘기집 원리의 분해
  증명은 지수 길이(Haken 1985). 모든 체계에 대한 초다항 하한은
  \(\mathrm{NP}\ne\mathrm{coNP}\)와 동치라 미해결. 괴델의 속도 향상
  정리로 ``가장 짧은 증명''은 체계에 의존한다.
\item
  \textbf{증명 동일성:} 정규화 기준(Prawitz), 일반성 기준(Lambek),
  증명망, Hughes의 조합적 증명. Straßburger(2019)가 증명 동일성의 네
  가지 목표 제시.
\item
  \textbf{2019년 Phil. Trans. R. Soc. A 특집호}, 자동 정리
  증명(Thiele--Wos 2002), 역수학.
\end{itemize}

단순성의 합의된 정의는 아직 없다.


\section{Lean 형식화}

2026-08-21.

Lean은 정리 증명기이자 함수형 프로그래밍 언어다. 정의·정리·증명을 기계가
한 단계씩 검사할 수 있는 형태로 다시 쓰는 것이 형식화다.

\begin{itemize}
\tightlist
\item
  기반: 귀납적 구성의 계산(Calculus of Inductive Constructions). 명제는
  타입, 증명은 그 타입의 항(커리--하워드 대응).
\item
  사례: Mathlib(공동 수학 라이브러리), Liquid Tensor Experiment, PFR
  추측의 형식화.
\item
  CS 쪽 연결: 의존 타입, 증명 보조기, 형식 검증(CompCert 컴파일러, seL4
  커널).
\end{itemize}

\subsection{직접 한 것}

\begin{itemize}
\tightlist
\item
  원 밖의 점에서 그은 접선 공식 \(r(px+qy)\pm\ell(py-qx)=rd^2\)를 Lean 4
  (v4.23.0, core only, Mathlib 없이)로 임의의 가환환 위에서 검증.
  회전·복소수 해석도 별도 파일로 검증. → \url{analytic-geometry.md}
\end{itemize}

\subsection{파일}

\begin{itemize}
\tightlist
\item
  \href{../lean/Tangent.lean}{\texttt{lean/Tangent.lean}}: 법선벡터
  \(n=(rp-s\ell q,\ rq+s\ell p)\)로 쓴 공식과 같음, 직선이 \(P\)를 지남,
  \(|n|^2=d^4\)(거리 \(=r\)), 원과의 공통점은 \(T=rn/d^2\) 하나뿐, 두
  접선(\(s=\pm1\))이 다름
\item
  \href{../lean/TangentSimple.lean}{\texttt{lean/TangentSimple.lean}}:
  \(P\)를 각 \(\alpha\)(\(\cos\alpha=r/d\), \(\sin\alpha=\ell/d\))만큼
  돌린 회전 형태, 복소수 형태 \(T\bar P=r(r+is\ell)\)
\end{itemize}

모두 Mathlib 없이 \texttt{grind} 전술로 증명하고, 정수 예제는
\texttt{decide}로 확인한다. 실행: Lean 4 v4.23.0에서
\texttt{lean\ lean/Tangent.lean}.


\end{document}
